%% ACRONYMES - Système ERP et Technologies

% Ressources Humaines et Gestion
\newacronym{RH}{RH}{Ressources Humaines}
\newacronym{ERP}{ERP}{Enterprise Resource Planning (Progiciel de Gestion Intégré)}
\newacronym{CV}{CV}{Curriculum Vitae}

% Intelligence Artificielle
\newacronym{IA}{IA}{Intelligence Artificielle}
\newacronym{LLM}{LLM}{Large Language Model (Modèle de Langage de Grande Taille)}

% Technologies Backend
\newacronym{ASGI}{ASGI}{Asynchronous Server Gateway Interface}
\newacronym{WSGI}{WSGI}{Web Server Gateway Interface}

% Protocoles et Authentification
\newacronym{REST}{REST}{Representational State Transfer}
\newacronym{JWT}{JWT}{JSON Web Token}
\newacronym{OAuth2}{OAuth2}{OAuth 2.0 - Protocole d'Autorisation}
\newacronym{OpenAPI}{OpenAPI}{OpenAPI Specification}

% Base de Données et Infrastructure
\newacronym{ORM}{ORM}{Object-Relational Mapping}
\newacronym{ACID}{ACID}{Atomicité, Cohérence, Isolation, Durabilité}
\newacronym{PostgreSQL}{PostgreSQL}{Système de Gestion de Base de Données Relationnelle}
\newacronym{Redis}{Redis}{Magasin de Données Clé-Valeur en Mémoire}
\newacronym{IoT}{IoT}{Internet des Objets}
\newacronym{PubSub}{Pub/Sub}{Publisher/Subscriber (Publication/Souscription)}

% Frontend et Frameworks
\newacronym{React}{React}{Bibliothèque JavaScript pour les Interfaces Utilisateur}
\newacronym{Vite}{Vite}{Outil de Compilation et Serveur de Développement}
\newacronym{SPA}{SPA}{Single Page Application}

% Outils et Librairies Nommés
\newacronym{FastAPI}{FastAPI}{Framework Python Web Asynchrone}
\newacronym{SQLAlchemy}{SQLAlchemy}{Toolkit SQL et Object-Relational Mapping pour Python}
\newacronym{Pydantic}{Pydantic}{Bibliothèque de Validation de Données Python}
\newacronym{Alembic}{Alembic}{Outil de Migration de Base de Données pour SQLAlchemy}
\newacronym{LangChain}{LangChain}{Framework pour Développer des Applications avec des LLM}
\newacronym{Docker}{Docker}{Plateforme de Conteneurisation}
\newacronym{DockerCompose}{Docker Compose}{Outil d'Orchestration Multi-Conteneurs pour le Développement}
\newacronym{Kubernetes}{Kubernetes}{Système d'Orchestration de Conteneurs}
\newacronym{OpenAI}{OpenAI}{Entreprise en Intelligence Artificielle et Créatrice de GPT}

% Entreprise
\newacronym{Harmony}{Harmony}{Harmony - Société Technologique Marocaine}

%% GLOSSAIRE - Termes et Concepts

\newglossaryentry{Streaming}
 {
 name=Streaming de Réponses,
 description={Transmission progressive des résultats au fur et à mesure de leur génération, sans attendre la réponse complète}
 }

\newglossaryentry{Kanban}
 {
 name=Kanban,
 description={Système de gestion de projets utilisant des colonnes (À faire, En cours, Fait) et des cartes pour canaliser le flux de travail}
 }

\newglossaryentry{Scrum}
 {
 name=Scrum,
 description={Framework agile pour la gestion de projets basé sur des sprints itératifs et une collaboration quotidienne}
 }

\newglossaryentry{WebSocket}
 {
 name=WebSocket,
 description={Protocole permettant une communication bidirectionnelle en temps réel entre un client et un serveur}
 }

\newglossaryentry{Notification}
 {
 name=Notification,
 description={Message d'alerte envoyé instantanément à un utilisateur pour l'informer d'un événement (changement de statut, deadline, etc.)}
 }
 